\section{Comparison tensors and the CM implication}\label{sec:hinge}

Put $k=\Qbar$.  On a smooth algebraic elliptic curve choose a
holomorphic de Rham class $\Omega$ and a complementary class $\Xi$,
normalized by the polarization.  Integration over an oriented integral
basis gives
\[
 M=\begin{pmatrix}\omega&h\\\omega_2&h_2\end{pmatrix},
 \qquad \det M=2\pi i,\qquad \tau=\omega_2/\omega\in\mathfrak H.
\]
For the standard families these are the classes and cycles of
Section~\ref{sec:families}.  Their period normalization turns an
identity into
\begin{equation}\label{eq:quadratic}
 a\omega^2+b\omega h=\alpha\pi,\qquad
 a=\tfrac32(A+Bu_s),\quad b=\tfrac32Bv_s,
 \quad(a,b)\ne(0,0).
\end{equation}
The nonvanishing theorem already proved gives $\alpha\ne0$.
This section identifies the precise comparison assertion needed to
deduce CM, and explains two equivalent ways to complete the information
in \eqref{eq:quadratic}.

\subsection{Quadratic genericity}
Let $\mathcal P_E$ be the torsor of tensor isomorphisms between the
de Rham and Betti fiber functors on the de Rham--Betti tensor category
generated by $H^1(E)$ and the Tate object, with algebraic coefficients.
At a non-CM elliptic curve its group is
$\operatorname{GL}_2$; after choosing the preceding bases its coordinate
ring is $k[w,h,v,k_2,(wk_2-vh)^{-1}]$.  The determinant is the comparison
coordinate of the Tate object.  Let $Z_E$ be the Zariski closure over
$k$ of the comparison point $M$, and set
\[
 U_\gamma=\langle w^2,wh,wk_2-vh\rangle_k\subset\mathcal O(\mathcal P_E).
\]
The three polynomials defining this space are linearly independent.
Evaluation factors through the actual closure:
\[
\begin{tikzcd}[column sep=large,row sep=large]
 U_\gamma\arrow[r,hook]\arrow[dr,"\mathrm{restriction}"']
 &\mathcal O(\mathcal P_E)\arrow[d,two heads]\arrow[r,"\operatorname{ev}_M"]
 &\mathbb C\\
 &\mathcal O(Z_E)\arrow[ur,hook]&
\end{tikzcd}
\]
The lower injection is induced by evaluation at $M$.  Thus
Hypothesis~\ref{hyp:QPI} is exactly injectivity of evaluation on
$U_\gamma$ for the specified non-CM fibers.  Full genericity of the
comparison point would assert injectivity on the entire coordinate ring.

\begin{proposition}\label{prop:CMforcing}
Hypothesis~\ref{hyp:QPI} forces every standard identity to have a CM fiber.
\end{proposition}
\begin{proof}
The positive endpoint has already been excluded.  If a remaining fiber
were non-CM, \eqref{eq:quadratic} would make the nonzero element
\[
 aw^2+bwh-\frac{\alpha}{2i}(wk_2-vh)\in U_\gamma
\]
vanish at $M$, contrary to injectivity.
\end{proof}

All results below concern either explicit numerical conjectures or
complete comparison tensors.  None identifies a functional
transcendence theorem with its numerical specialization, nor infers
the remaining coordinates of a tensor from one algebraic coordinate.
The equivalent restricted rational-direction formulation and the
obstructions to such inferences are collected in
Appendix~\ref{rev:app-comparison}.

\subsection{The modular two-jet}\label{subsec:modular-jets}
For this subsection take the Weierstrass model
$y^2=4X^3-g_2X-g_3$, with $\Omega=dX/y$ and $\Xi=X\,dX/y$.
Write
\[
 j_0=j(\tau),\qquad P=\frac{\omega^2}{\pi},\qquad
 Q=\frac{\omega h}{\pi},\qquad
 R(T)=\frac{2}{3T}+\frac1{2(T-1728)}.
\]
Primes on $j$ denote ordinary $\tau$-derivatives.

\begin{proposition}[Period--jet comparison]\label{prop:period-jets}
If $j_0\notin\{0,1728\}$, then $j'(\tau)\ne0$ and
\begin{align}
 j'&=-9ij_0\frac{g_3}{g_2}P,\label{eq:jet-first}\\
 \frac{j''}{j'}&=R(j_0)j'-iQ.\label{eq:jet-second}
\end{align}
Consequently quadratic-period independence is equivalent to either
of the following linear independence assertions over $k$:
\[
 \{1,j',j''/j'\},\qquad \{j',j'^2,j''\}.
\]
Moreover,
\begin{equation}\label{eq:hierarchy-jets-field}
 k(P,Q)=k(j',j'').
\end{equation}
\end{proposition}
\begin{proof}
With $D=(2\pi i)^{-1}d/d\tau$, the lattice formulas in
\cite[Proposition~4.7]{Fonseca}, in the convention $h=-\eta$, are
\begin{equation}\label{eq:hierarchy-eisenstein}
 E_2=-\frac{3Q}{\pi},\qquad
 E_4=\frac{3g_2P^2}{4\pi^2},\qquad
 E_6=\frac{27g_3P^3}{8\pi^3}.
\end{equation}
Ramanujan's differential equations
\[
 DE_2=(E_2^2-E_4)/12,\quad
 DE_4=(E_2E_4-E_6)/3,\quad
 DE_6=(E_2E_6-E_4^2)/2
\]
applied to $j=1728E_4^3/(E_4^3-E_6^2)$ give
\[
 Dj=-jE_6/E_4,\qquad
 \frac{j''}{2\pi i j'}=\frac{E_2}{6}
       -\frac{2E_6}{3E_4}-\frac{E_4^2}{2E_6}.
\]
Substitute \eqref{eq:hierarchy-eisenstein} and
$g_2^3/g_3^2=27j_0/(j_0-1728)$.  This proves
\eqref{eq:jet-first} and \eqref{eq:jet-second}; every denominator is nonzero.  The change from
$(P,Q)$ to $(j',j''/j')$ is invertible and $k$-linear.  Multiplication
by $j'$ proves the second independence equivalence, and the same
formulas and their inverses prove the field equality.
\end{proof}

\begin{corollary}\label{cor:forbidden-jet}
The relation \eqref{eq:quadratic} is equivalent to
\begin{equation}\label{eq:forbidden-jet}
 b j''+\left(\frac{ag_2}{9j_0g_3}-bR(j_0)\right)j'^2
                  +i\alpha j'=0.
\end{equation}
Unconditionally, $j''/j'^2$ is transcendental, and the relation space
of $(j',j'^2,j'')$ over $k$ has dimension at most one.
\end{corollary}
\begin{proof}
The inverse dictionary is
$P=ig_2j'/(9j_0g_3)$ and
$Q=i(j''/j'-R(j_0)j')$.  Substitution proves the first assertion.
It also gives
\[
 \frac{j''}{j'^2}=R(j_0)+\frac{g_2}{9j_0g_3}\frac h\omega.
\]
Chudnovsky's theorem, in the form used in the preceding section,
makes $h/\omega$ transcendental.  Thus $j'^2$ and $j''$ are linearly
independent; adjoining one further number creates at most one relation.
\end{proof}

For $b=0$, the missing implication already contains the restriction
of Lang's derivative conjecture to these points: algebraic $j(\tau)$
and $j'(\tau)\ne0$ should imply transcendental $j'(\tau)$
\cite[p.~167, Conjecture~34]{WaldschmidtSurvey}.  For $b\ne0$ it is
the affine quadratic obstruction \eqref{eq:forbidden-jet}.

\subsection{Sufficient numerical conjectures}
\label{subsec:precise-period-hierarchy}
Let
\[
 K_E=k(\omega,\omega_2,h,h_2),\qquad q=e^{2\pi i\tau}.
\]
At a non-CM algebraic elliptic curve, the following numerical forms of
the cited conjectures apply:
\begin{enumerate}
\item The elliptic period conjecture asserts
$\operatorname{trdeg}_k K_E=4$
\cite[Conjecture~4.5]{Fonseca}.
\item Fonseca's modular-value conjecture asserts
\[
 \operatorname{trdeg}_k
 k(2\pi i,\tau,q,E_2,E_4,E_6)\ge5
\]
at nonquadratic $\tau$ \cite[Conjecture~4.8]{Fonseca}.
\item The one-point modular Schanuel conjecture with derivatives asserts
\[
 \operatorname{trdeg}_{\Q}
 \Q(\tau,j(\tau),j'(\tau),j''(\tau))\ge3
\]
at nonquadratic $\tau$
\cite[Conjecture~2.3]{EterovicMSCD}.  This is the one-point case of the
lower bound three times the number of distinct
$\operatorname{GL}_2^+(\Q)$-orbits outside the imaginary quadratic locus.
\end{enumerate}
Only their restrictions to the algebraic $j$-values in question are needed.
Denote the second numerical bound by $\mathrm{FV}_\tau$ and the third
by $\mathrm{MSCD}_\tau$.

\begin{proposition}[Hierarchy of sufficient inputs]
\label{prop:hierarchy-implications}\label{prop:mscd-cm}\label{prop:hierarchy-fields}
At each such non-CM fiber,
\[
\begin{tikzcd}[column sep=small,row sep=large]
 \mathrm{FV}_\tau
 \arrow[r,Rightarrow]
 &\operatorname{trdeg}_k K_E=4\arrow[d,Rightarrow]&\\
 \mathrm{MSCD}_\tau\arrow[r,Rightarrow]
 &\operatorname{trdeg}_k k(\tau,j',j'')=3
 \arrow[r,Rightarrow]
 &\text{QPI}.
\end{tikzcd}
\]
Each of these conjectural inputs therefore suffices for the
CM implication and the conditional classification.
\end{proposition}
\begin{proof}
Legendre's identity and \eqref{eq:hierarchy-jets-field} give
\begin{equation}\label{eq:hierarchy-legendre}
 h_2=\tau h+\frac{2\pi i}{\omega},\qquad
 K_E=k(\omega,h,\tau,\pi),\qquad
 [K_E:k(\pi,\tau,j',j'')]\le2.
\end{equation}
Indeed, adjoining $\omega$ with $\omega^2=\pi P$ suffices.
Furthermore \eqref{eq:hierarchy-eisenstein} and its inverses
\[
 P=\frac{2g_2\pi}{9g_3}\frac{E_6}{E_4},\qquad
 Q=-\frac{\pi E_2}{3}
\]
give the exact equality
\begin{equation}\label{eq:hierarchy-modular-field}
 k(2\pi i,\tau,q,E_2,E_4,E_6)=k(\pi,\tau,q,P,Q).
\end{equation}
Thus the modular-value bound implies
$\operatorname{trdeg}_k K_E(q)\ge5$.  Adjoining $q$ raises
transcendence degree by at most one, whereas $K_E$ has four generators;
hence $\operatorname{trdeg}_k K_E=4$.

The finite extension in \eqref{eq:hierarchy-legendre} then makes
$\pi,\tau,j',j''$ algebraically independent.  Deleting $\pi$ gives
the middle assertion; MSCD gives it directly since $j(\tau)\in k$.
Deleting $\tau$ makes $j',j''$ algebraically independent, so the
nonzero polynomial \eqref{eq:forbidden-jet} cannot vanish.  Equivalently,
\eqref{eq:hierarchy-jets-field} makes $1,P,Q$ linearly independent.
\end{proof}

These implications do not assert equivalence of the named conjectures.
In fact the modular-value bound here is equivalent to the elliptic
period conjecture together with transcendence of $q$ over $K_E$.
The exact equalities above prove this by counting transcendence degree.
MSCD requires neither this nome extension nor the additional
independence of $\pi$.

\subsection{Completing the comparison tensor}\label{sec:comparison-coordinate}
\label{subsec:ksv-completion}
Write $H=H^1(E)$ and use ordinary symmetric powers, with the usual
sign modification of cohomological symmetry.  Polarization identifies
$\det H$ with the Tate object of period $2\pi i$.  There is a canonical
isomorphism
\[
 \Phi:\operatorname{Sym}^2H\otimes(\det H)^{-1}
       \xrightarrow{\sim}\operatorname{End}_0(H).
\]
For a two-dimensional realization $V$, it is defined by
\[
 \Phi(vw\otimes\delta)(x)
   =\tfrac12\{\delta(w\wedge x)v+\delta(v\wedge x)w\},
 \qquad \delta\in(\det V)^\vee.
\]
This definition is invariant under change of basis.  It proves, in
particular, the commutativity of the comparison square
\[
\begin{tikzcd}[column sep=large,row sep=large]
 \bigl(\operatorname{Sym}^2H_{\mathrm{dR}}\otimes
       (\det H_{\mathrm{dR}})^{-1}\bigr)_{\mathbb C}
 \arrow[r,"\Phi_{\mathrm{dR}}"]\arrow[d,"\operatorname{comp}"']
 &\operatorname{End}_0(H_{\mathrm{dR},\mathbb C})
 \arrow[d,"T\mapsto MTM^{-1}"]\\
 \bigl(\operatorname{Sym}^2H_{\mathrm B}\otimes
       (\det H_{\mathrm B})^{-1}\bigr)_{\mathbb C}
 \arrow[r,"\Phi_{\mathrm B}"]
 &\operatorname{End}_0(H_{\mathrm B,\mathbb C}).
\end{tikzcd}
\]
Here a subscript $\mathbb C$ denotes extension of scalars from $k$
or $\Q$, respectively.  A \emph{complete algebraic de Rham--Betti
tensor} is an algebraic de Rham tensor whose entire comparison image
has algebraic Betti coordinates.

The exact external input is as follows.
\begin{theorem}[Kreutz--Shen--Vial, specialized]
\label{thm:ksv-specialized}
For one elliptic curve over $k$ and a coefficient field $L\subseteq k$,
comparison-compatible endomorphisms with algebraic de Rham matrix and
Betti matrix over $L$ form $\operatorname{End}(E)\otimes_{\Z}L$.
With algebraic coefficients, the de Rham--Betti group generated by
$H$ is $\operatorname{GL}_2$ in the non-CM case and the torus of the
CM field in the CM case.  Every complete algebraic de Rham--Betti tensor
on a power of $E$, with a Tate twist, is an algebraic linear combination
of algebraic cycle classes.
\end{theorem}
The endomorphism assertion is
\cite[Theorem~5.10 and Corollary~5.11]{KSV}; the tensor and group
assertions are \cite[Theorem~5.13]{KSV}.  The latter theorem allows
algebraic coefficient fields containing at most one of the distinct
CM fields among the factors.  A power of one elliptic curve satisfies
this condition, including for $L=k$.

\begin{proposition}[Exact completion]\label{prop:ksv-exact-completion}
Assume \eqref{eq:quadratic} and $b\ne0$.  Put
$\nu=a\Omega+b\Xi$ and $c=2ib/\alpha$.  Let $\Pi$ be the de Rham
projector onto $k\nu$ along $k\Omega$.  Then the following are equivalent:
\begin{enumerate}
\item $E$ has CM;
\item $\Pi$ is a complete algebraic de Rham--Betti endomorphism;
\item the tensor $t=\Omega\nu/(\Omega\wedge\nu)$ is complete;
\item any one of the three Betti entries of $\Pi$ other than its
upper-right entry is algebraic.
\end{enumerate}
At CM, its Betti entries generate precisely the CM field, and
\begin{equation}\label{eq:ksv-cm-field}
 c=\overline\tau-\tau,\qquad
 b=-\alpha\operatorname{Im}\tau,\qquad
 (\alpha/b)^2\in\Q_{>0}.
\end{equation}
\end{proposition}
\begin{proof}
Put $u=a\omega+bh$.  The comparison matrix in the basis $(\Omega,\nu)$
has columns $(\omega,\tau\omega)^t$ and
$(u,(\tau+c)u)^t$: its determinant is $2\pi ib$, whereas
$\omega u=\alpha\pi$.  Conjugating $\operatorname{diag}(0,1)$ gives
\begin{equation}\label{eq:ksv-projector}
 \Pi_{\mathrm B}=\frac1c
 \begin{pmatrix}-\tau&1\\-\tau(\tau+c)&\tau+c\end{pmatrix}.
\end{equation}
Algebraicity of either diagonal entry makes $\tau$ algebraic;
algebraicity of the lower-left entry gives a quadratic equation for
$\tau$.  Conversely, algebraic $\tau$ makes every entry algebraic.
Schneider's theorem identifies algebraicity of $\tau$, since
$j(E)\in k$, with CM.  Alternatively, a complete rank-one projector
cannot belong to the non-CM endomorphism algebra $k$, by
Theorem~\ref{thm:ksv-specialized}.

The defining formula for $\Phi$ gives
$\Phi(t)=\Pi-\tfrac12I$, proving the equivalence with tensor completion.
At CM the algebra $\operatorname{End}(E)\otimes k$ splits as
$k\times k$.  Its two rank-one projectors select the CM eigenspaces.
Since $\Pi$ kills the holomorphic eigenspace of ratio $\tau$, its image
has ratio $\overline\tau$.  Thus $\tau+c=\overline\tau$, proving the
first two identities in \eqref{eq:ksv-cm-field}.  Its entries generate
$\Q(c,\tau)=\Q(\tau)$: recover $c$ from the upper-right entry and
$\tau$ from the upper-left entry.  For imaginary quadratic $\tau$,
$(\operatorname{Im}\tau)^2$ is positive rational, proving the last
assertion.
\end{proof}

This argument exhibits the exact missing datum.  The scalar identity
fixes $\Pi_{12}=1/c$, while completion needs one more algebraic entry.
For $b=0$ the associated tensor is instead a nonzero nilpotent in
$\operatorname{End}_0(H)$; it can never be complete, since neither
$k$ nor $k\times k$ has a nonzero nilpotent.  This identifies the
different algebraic type of the constant-weight obstruction.

\subsection{Propagation to a second cycle}\label{sec:two-cycles}
\begin{proposition}\label{prop:two-cycles}
Under \eqref{eq:quadratic}, $E$ has CM if and only if the same
coefficients $a,b$ give a value in $k\pi$ on some integral cycle
$m\gamma_1+n\gamma_2$ with $n\ne0$.
\end{proposition}
\begin{proof}
Set $L=m+n\tau$.  Legendre's identity gives
$h_{m,n}=Lh+2\pi in/\omega$ and $\omega_{m,n}=L\omega$.  Hence
\begin{equation}\label{eq:cycle-propagation}
 \frac{\omega_{m,n}(a\omega_{m,n}+bh_{m,n})}{\pi}
       =\alpha L^2+2ibnL.
\end{equation}
Algebraicity of this expression, with $n\ne0$, gives a quadratic
equation for $\tau$ whose leading coefficient is $\alpha n^2\ne0$.
Schneider's theorem proves CM.  Conversely, CM makes $\tau$ algebraic,
so the expression belongs to $k$ on every integral cycle.
\end{proof}
The second cycle expresses geometrically what the omitted comparison
coordinate expresses algebraically.  The original identity supplies
neither completion nor propagation; each is an exact additional target
for an unconditional CM proof.
